\scanpage{24}
% L1–2
\uline{Corollaire} Les conditions (i) (ii) impliquent que si $u\colon X\to Y$ est \uline{surjectif}, alors $F(Y)\to F(X)$ est \uline{injectif}.

% L3–6
En effet, grâce à (i) on peut supposer \deleted{\unsure{un|que}} $X$, $Y$ réduits, Mais alors $u$ est un épim, et nous avons déjà vu que \deleted{\unsure{on a}} alors $F(Y)\mathrel{-}F(X)$ est injectif. OK.

% L4–6, left
\diagram{\begin{tikzcd}[ampersand replacement=\&,column sep=small,row sep=small]
F(X)\arrow[d] \& F(Y)\arrow[l]\arrow[d,"\mathrm{inj}"]\\
F(X_{\mathrm{réd}}) \& F(Y_{\mathrm{réd}})\arrow[l]
\end{tikzcd}}

% L7
Cas particulier

% L8–12
\uline{Proposition} \deleted{\unsure{Soient}} Soit $A$ préschéma abélien sur $S$ loc noeth, et pour tt $S'$ sur $S$, soit $F(S')$ l'ens des sous-préschémas abéliens de $A'=A\unsure{\otimes|\times}_S\unsure{S'_0|S'}$. Alors si $S'$ est tel que $S'\mathrel{-}S$ soit un épimorphisme, on a une suite exacte (où $S''=S'\times_S S'$)

% L13
\[
0\longrightarrow F(S)\longrightarrow F(S')\rightrightarrows F(S'')
\]

% L14–16
En effet, la condition (i) du Lemme résulte de la rigidité, la condition (ii) résulte de la théorie de la descente.

% L17–19
\deleted{\uline{Proposition}} \inserted{Corollaire} \inserted{Énoncé analogue pour les hom} \deleted{(\unsure{schéma|Schéma} \unsure{Krizien|Krizmen}) Supposons que $S$ soit le spectre d'} de schémas abéliens.

% archive sheet number: 23

